% !TEX root = ../main.tex
\chapter{Kết luận và hướng phát triển}

\section{Mức độ hoàn thành mục tiêu}

Đồ án đã xây dựng RabbitCV với trọng tâm là tạo CV và hỗ trợ nội dung bằng AI; các chức năng tuyển dụng đóng vai trò cung cấp ngữ cảnh sử dụng CV. Mức độ hoàn thành được tổng hợp từ thành phần triển khai và kết quả đánh giá ở Chương 5, không chỉ từ việc giao diện có thể hiển thị.

\begin{table}[H]
    \centering
    \small
    \begin{tabularx}{\textwidth}{|p{0.24\textwidth}|Y|p{0.17\textwidth}|} \hline
        \textbf{Mục tiêu} & \textbf{Kết quả và bằng chứng} & \textbf{Kết luận} \\ \hline
        Tạo và quản lý CV & Resume có cấu trúc, template, preview A4, in/tải PDF; test A4 và ngày tháng đạt. & Hoàn thành trong phạm vi \\ \hline
        AI hỗ trợ CV & Viết/cải thiện tiểu sử, review, fill, so khớp CV--Job và phỏng vấn; có sanitizer, schema, quota và usage log. & Hoàn thành mức demo \\ \hline
        Tìm việc và ứng tuyển & Tìm/lưu Job, nộp Resume hoặc PDF, unique index, trạng thái và jobSnapshot. & Hoàn thành trong phạm vi \\ \hline
        Kết nối hai phía & Cổng doanh nghiệp, xử lý Application, chat Socket.IO và notification. & Hoàn thành một phần \\ \hline
        Quản trị & Có thống kê và nhóm endpoint quản lý User, Company, Job. & Hoàn thành một phần \\ \hline
        Kiểm chứng kỹ thuật & 41/41 test đạt, 0 lỗi lint, build thành công; còn hai cảnh báo hook và cảnh báo bundle. & Đạt baseline \\ \hline
    \end{tabularx}
    \caption{Tổng hợp mức độ hoàn thành mục tiêu}
    \label{tab:objective-completion}
\end{table}

Kết quả nổi bật của đồ án không nằm ở số lượng màn hình mà ở việc nối được vòng đời CV với nghiệp vụ sử dụng: dữ liệu CV được chỉnh sửa có cấu trúc, được trình bày theo A4, có thể nhận gợi ý AI dưới quyền quyết định của người dùng, rồi được dùng để ứng tuyển và trao đổi với doanh nghiệp.

\section{Đóng góp kỹ thuật}

Thứ nhất, hệ thống tích hợp AI theo một đường biên rõ ràng. Frontend không giữ khóa provider; backend xác thực, giới hạn tần suất và quota, làm sạch dữ liệu, yêu cầu JSON Schema và kiểm tra kết quả. Cơ chế ``gợi ý rồi mới áp dụng'' tránh để mô hình tự sửa hồ sơ cá nhân.

Thứ hai, một số tình huống dễ gây mất hoặc lộ dữ liệu đã được xử lý ngay trong thiết kế: \texttt{jobSnapshot} giữ ngữ cảnh lịch sử; unique index ngăn ứng tuyển trùng; Resume ảo tách PDF ứng tuyển khỏi thư viện CV; truy vấn danh sách không trả \texttt{pdfData}; trạng thái đọc và ẩn chat được lưu theo từng người dùng.

Thứ ba, báo cáo gắn yêu cầu với mã nguồn và test. Việc công bố cả 41 test đạt lẫn phần chưa đánh giá --- browser end-to-end, database integration, phân quyền đầy đủ và chất lượng ngữ nghĩa AI --- giúp giới hạn kết luận đúng với bằng chứng hiện có.

\section{Bài học kinh nghiệm}

\begin{itemize}
    \item Dữ liệu nghề nghiệp cần tách nội dung cấu trúc khỏi bản trình bày. Nếu chỉ lưu PDF, việc đổi template, review theo từng mục và hỗ trợ AI sẽ khó kiểm soát.
    \item Xóa hoặc ẩn dữ liệu trong hệ thống nhiều bên phải xét quyền sở hữu và nhu cầu lưu lịch sử. Snapshot và trạng thái ẩn theo user phù hợp hơn xóa dây chuyền.
    \item Kiểm tra ở frontend không thay thế kiểm tra server. Role, userId, MIME hoặc trạng thái Job từ trình duyệt đều phải được xem là dữ liệu không đáng tin.
    \item Structured output làm giảm lỗi tích hợp AI nhưng không bảo đảm nội dung đúng. Chất lượng mô hình cần rubric, bộ dữ liệu cố định và người chấm độc lập.
    \item Một PDF ``trông đúng'' trên màn hình chưa đủ; cần kiểm tra kích thước vật lý, ngắt trang, Print Preview và trình xem nhiều trang.
\end{itemize}

\section{Hạn chế hiện tại}

Hệ thống chưa được xem là sẵn sàng production. PDF base64 làm document MongoDB lớn; backend còn tập trung nhiều route; notification dùng polling; cron chưa có khóa phân tán; một số màn hình admin chưa đồng bộ đủ API. Bộ test chưa chạy MongoDB thật, trình duyệt end-to-end, tải, pentest hoặc mạng Socket.IO không ổn định. Working tree sản phẩm tại thời điểm đánh giá cũng chưa được đóng thành commit/tag tái lập. Cuối cùng, hiệu quả ngữ nghĩa và tính công bằng của AI chưa được đánh giá độc lập.

\section{Lộ trình phát triển}

\begin{table}[H]
    \centering
    \small
    \begin{tabularx}{\textwidth}{|p{0.18\textwidth}|Y|} \hline
        \textbf{Giai đoạn} & \textbf{Công việc ưu tiên} \\ \hline
        Ngắn hạn & Commit/tag baseline DA2; sửa hai cảnh báo React Hook; hoàn thiện API admin; thêm integration test cho Application unique/snapshot và ma trận quyền; chạy checklist nghiệm thu trên trình duyệt. \\ \hline
        Trung hạn & Chuyển PDF/ảnh sang object storage hoặc GridFS; tách backend theo domain controller--service; đưa notification quan trọng lên Socket.IO; thêm queue/worker cho tác vụ file và cron; bổ sung log/audit. \\ \hline
        Dài hạn & Xây dựng bộ dữ liệu chấm AI có người đánh giá; theo dõi chi phí/chất lượng theo model; kiểm thử accessibility, tải và bảo mật; thiết kế quy trình backup, khôi phục và quan sát production. \\ \hline
    \end{tabularx}
    \caption{Lộ trình phát triển đề xuất}
    \label{tab:roadmap}
\end{table}

\section{Tổng kết}

RabbitCV chứng minh tính khả thi của một ứng dụng web hỗ trợ viết CV có tích hợp AI, đồng thời cho thấy AI chỉ là một thành phần trong hệ thống dữ liệu, quyền truy cập và nghiệp vụ rộng hơn. Sản phẩm hiện đáp ứng mục tiêu đồ án ở mức nguyên mẫu có thể kiểm chứng: chức năng cốt lõi đã hiện thực, các biện pháp kiểm soát quan trọng đã có và quy trình tự động đã vượt qua toàn bộ 41 test. Những giới hạn còn lại đã được xác định thành công việc cụ thể cho giai đoạn hoàn thiện và triển khai thực tế.
